\section{Introduction}

A time--memory tradeoff (TMTO) moves work from the online phase of a search
problem into a reusable preprocessing phase.  In Hellman's classical
construction~\cite{Hellman1980}, a large key or state space is represented by
iteration chains while only chain endpoints are retained.  The construction
therefore occupies a middle ground between exhaustive search, which uses little
memory and large online time, and complete table lookup, which uses large
memory and very small online time.

The central implementation issue is that a chain table does not contain
$m t$ independent samples.  Once two iterates of the same mapping meet, their
future trajectories coincide; a chain can also revisit one of its own previous
states.  Consequently, the effective coverage of a table is the number of
\emph{distinct} represented states, not its nominal number of positions.
This phenomenon is part of the classical TMTO literature rather than a new
observation: Hellman's analysis, later work on random mappings, success
probability, perfect tables, false alarms, and rainbow reductions all treat
different aspects of collisions, merging, and coverage
~\cite{Hellman1980,FlajoletOdlyzko1990,Oechslin2003,AvoineJunodOechslin2008,MaHong2009,Hong2008}.

The purpose of this paper is therefore narrower and experimentally explicit.
We ask how accurately established random-function TMTO theory describes a
\emph{concrete finite cryptographic mapping} when every represented state can
be enumerated exactly, and whether independent implementations reproduce the
same semantics and finite-state statistics.

\paragraph{Concrete iteration map.}
For a reduced state $x$, a canonical DES-key encoding $K(x)$, a fixed plaintext
$P$, encryption $E$, and reduction $R$, we study
\[
    F(x)=R(E_{K(x)}(P)).
\]
DES is obsolete and is used only as a compact deterministic fixture.  The state
dimension is deliberately reduced so that exact coverage, cycles, endpoint
multiplicities, and recovery behavior can be computed rather than estimated
from sampling.

\paragraph{A motivating discrepancy and its correct baseline.}
The repository's initial $N=2^{16}$ experiment used $m=1024$ fixed-length
chains of length $t=64$, so $mt=N$.  Treating the $mt$ chain positions as
independent samples would give the occupancy expression
\[
    1-e^{-mt/N}=1-e^{-1}\approx 63.2\%.
\]
That number is \emph{not} the expected coverage of a Hellman matrix.  Using the
random-function recurrence of Ma and Hong~\cite{MaHong2009}, the same
parameters predict approximately $16.651\%$ coverage of the state space.  The
measured median exact coverage is approximately $16.8\%$.  Thus the initial
measurement is in close agreement with established random-function theory.

This distinction also exposes an important parameterization point.  The
standard Hellman matrix-stopping regime is expressed in terms of $m t^2$
relative to $N$, not merely $mt/N$.  Our initial configuration satisfies
$mt=N$ but $m t^2=64N$; it should therefore be treated as a finite-state stress
point, not as a standard matrix-stopping configuration.

\paragraph{Contributions and scope.}
This paper is an experimental and reproducibility study, not a claim to have
discovered chain coalescence or classical TMTO coverage theory.  Its intended
contributions are:
\begin{enumerate}
  \item a precise, executable reduced-DES state model with canonical DES parity
        handling and an explicit reduction;
  \item exact coverage and exact-recovery instrumentation for Hellman and
        bounded distinguished-point tables;
  \item direct comparison of concrete cryptographic mappings with established
        random-function coverage predictions over controlled finite-state
        parameter sweeps;
  \item independent Java and Rust reference implementations, with shared
        semantic fixtures and, before any timing comparison, a common committed
        experiment plan;
  \item a reproducibility artifact in which numerical claims are generated from
        committed inputs, outputs, scripts, and environment metadata.
\end{enumerate}

We make no claim to break full-size DES or AES, and no implementation in this
repository is intended for production cryptography.
